\documentclass[10pt,a4paper]{amsart}
\usepackage{amsmath,amssymb,amsthm,graphicx}
\usepackage{verbatim}
\newenvironment{reference}{\comment}{\endcomment}
%\newenvironment{reference}{}{}
\newenvironment{slogan}{\comment}{\endcomment}
\newenvironment{history}{\comment}{\endcomment}

% For commutative diagrams we use Xy-pic
\usepackage[all]{xy}

% We use 2cell for 2-commutative diagrams.
\xyoption{2cell}
% Theorem environments.
%
\theoremstyle{plain}
\newtheorem{theorem}[subsection]{Theorem}
\newtheorem{proposition}[subsection]{Proposition}
\newtheorem{lemma}[subsection]{Lemma}

\theoremstyle{definition}
\newtheorem{definition}[subsection]{Definition}
\newtheorem{example}[subsection]{Example}
\newtheorem{exercise}[subsection]{Exercise}
\newtheorem{situation}[subsection]{Situation}

\theoremstyle{remark}
\newtheorem{remark}[subsection]{Remark}
\newtheorem{remarks}[subsection]{Remarks}

\numberwithin{equation}{subsection}

% Macros
%
\def\lim{\mathop{\mathrm{lim}}\nolimits}
\def\colim{\mathop{\mathrm{colim}}\nolimits}
\def\Spec{\mathop{\mathrm{Spec}}}
\def\Hom{\mathop{\mathrm{Hom}}\nolimits}
\def\Ext{\mathop{\mathrm{Ext}}\nolimits}
\def\SheafHom{\mathop{\mathcal{H}\!\mathit{om}}\nolimits}
\def\SheafExt{\mathop{\mathcal{E}\!\mathit{xt}}\nolimits}
\def\Sch{\mathit{Sch}}
\def\Mor{\mathop{\mathrm{Mor}}\nolimits}
\def\Ob{\mathop{\mathrm{Ob}}\nolimits}
\def\Sh{\mathop{\mathit{Sh}}\nolimits}
\def\NL{\mathop{N\!L}\nolimits}
\def\CH{\mathop{\mathrm{CH}}\nolimits}
\def\proetale{{pro\text{-}\acute{e}tale}}
\def\etale{{\acute{e}tale}}
\def\QCoh{\mathit{QCoh}}
\def\Ker{\mathop{\mathrm{Ker}}}
\def\Im{\mathop{\mathrm{Im}}}
\def\Coker{\mathop{\mathrm{Coker}}}
\def\Coim{\mathop{\mathrm{Coim}}}

% Boxtimes
%
\DeclareMathSymbol{\boxtimes}{\mathbin}{AMSa}{"02}

%
% Macros for moduli stacks/spaces
%
\def\QCohstack{\mathcal{QC}\!\mathit{oh}}
\def\Cohstack{\mathcal{C}\!\mathit{oh}}
\def\Spacesstack{\mathcal{S}\!\mathit{paces}}
\def\Quotfunctor{\mathrm{Quot}}
\def\Hilbfunctor{\mathrm{Hilb}}
\def\Curvesstack{\mathcal{C}\!\mathit{urves}}
\def\Polarizedstack{\mathcal{P}\!\mathit{olarized}}
\def\Complexesstack{\mathcal{C}\!\mathit{omplexes}}
% \Pic is the operator that assigns to X its picard group, usage \Pic(X)
% \Picardstack_{X/B} denotes the Picard stack of X over B
% \Picardfunctor_{X/B} denotes the Picard functor of X over B
\def\Pic{\mathop{\mathrm{Pic}}\nolimits}
\def\Picardstack{\mathcal{P}\!\mathit{ic}}
\def\Picardfunctor{\mathrm{Pic}}
\def\Deformationcategory{\mathcal{D}\!\mathit{ef}}
\usepackage{cleveref}
\crefname{originalreference}{original reference}{original references}
\Crefname{originalreference}{Original reference}{Original references}
\makeatletter
\newlabel{lemma-chow-ring-projective-bundle}{{02TX}{source}{}{sourcepointer-1}{}}
\newlabel{lemma-chow-ring-projective-bundle@cref}{{[originalreference][0][]02TX}{[1][0][]source}{}{}{}}
\makeatother
\crefname{originalreference}{original reference}{original references}
\Crefname{originalreference}{Original reference}{Original references}
\makeatletter
\newlabel{lemma-linebundle}{{0B74}{source}{}{sourcepointer-2}{}}
\newlabel{lemma-linebundle@cref}{{[originalreference][0][]0B74}{[1][0][]source}{}{}{}}
\makeatother
\crefname{originalreference}{original reference}{original references}
\Crefname{originalreference}{Original reference}{Original references}
\makeatletter
\newlabel{lemma-pullback-affine-fibres-surjective}{{02TT}{source}{}{sourcepointer-3}{}}
\newlabel{lemma-pullback-affine-fibres-surjective@cref}{{[originalreference][0][]02TT}{[1][0][]source}{}{}{}}
\makeatother
\crefname{originalreference}{original reference}{original references}
\Crefname{originalreference}{Original reference}{Original references}
\makeatletter
\newlabel{lemma-pushforward-cap-c1}{{02SU}{source}{}{sourcepointer-4}{}}
\newlabel{lemma-pushforward-cap-c1@cref}{{[originalreference][0][]02SU}{[1][0][]source}{}{}{}}
\makeatother
\crefname{originalreference}{original reference}{original references}
\Crefname{originalreference}{Original reference}{Original references}
\makeatletter
\newlabel{lemma-relative-effective-cartier}{{02TR}{source}{}{sourcepointer-5}{}}
\newlabel{lemma-relative-effective-cartier@cref}{{[originalreference][0][]02TR}{[1][0][]source}{}{}{}}
\makeatother
\crefname{originalreference}{original reference}{original references}
\Crefname{originalreference}{Original reference}{Original references}
\makeatletter
\newlabel{lemma-restrict-to-open}{{02RX}{source}{}{sourcepointer-6}{}}
\newlabel{lemma-restrict-to-open@cref}{{[originalreference][0][]02RX}{[1][0][]source}{}{}{}}
\makeatother
\crefname{originalreference}{original reference}{original references}
\Crefname{originalreference}{Original reference}{Original references}
\makeatletter
\newlabel{situation-setup}{{02QL}{source}{}{sourcepointer-7}{}}
\newlabel{situation-setup@cref}{{[originalreference][0][]02QL}{[1][0][]source}{}{}{}}
\makeatother
\crefname{originallemma}{original lemma}{original lemmas}
\Crefname{originallemma}{Original lemma}{Original lemmas}
\makeatletter
\newlabel{spaces-divisors-lemma-characterize-OD}{{0847}{source}{}{sourcepointer-8}{}}
\newlabel{spaces-divisors-lemma-characterize-OD@cref}{{[originallemma][0][]0847}{[1][0][]source}{}{}{}}
\makeatother
\title{Pullback to the total space of a vector bundle is an isomorphism on Chow groups}
\author{The Stacks Project Authors}
\date{stacks-project:spaces-chow.tex:local-snapshot-20261012; GNU FDL 1.2 or later}
\begin{document}
\maketitle
\noindent\textit{Editorial note.} All delivered mathematical text is the author's own, copied from the pinned Stacks Project chapter file spaces-chow.tex; no statement or proof step was rewritten, reordered or removed.\par\medskip
\noindent\textit{Editorial note.} The counted proof is the author's complete argument: a completion to a projective bundle with the vector bundle realised as the complement of the zero section divisor, and an expansion through the projective bundle formula together with a vanishing pushforward computation\par\medskip
\noindent\textit{Editorial note.} The label lemma-vectorbundle names the counted statement, and the proof environment immediately following it in the source is delivered verbatim as its proof.\par\medskip
\noindent\textit{Editorial note.} The counted statement is carried by the author's own lemma environment in the chapter source.\par\medskip
\noindent\textit{Allowed standard prerequisites (editorial).}\begin{itemize}
\item Surjectivity of pullback for a flat morphism of relative dimension r with affine r-space fibres, applied to the vector bundle
\item The projective bundle formula for the projective bundle of the finite locally free sheaf with a trivial summand adjoined, and the identification of the complement of the zero section with the vector bundle
\item The exact sequence describing the kernel of restriction to an open subspace as the cycles supported on the complement, the relation between the invertible sheaf of an effective Cartier divisor and the tautological quotient, and the compatibility of cap with pushforward
\end{itemize}\medskip
\ifdefined\hypertarget\hypertarget{sourcepointer-1}{}\fi
\noindent\textit{Source reference.} Original reference 02TX: sources/stacks-spaces-chow/pinned-original-source.tex, source lines 4381--4381. Prerequisite pointer invoked at proof line(s) 4558; it names a standard fact of the same chapter that the counted argument uses, and every step of the counted proof built on it is retained verbatim in the delivered body.\par\medskip
\ifdefined\hypertarget\hypertarget{sourcepointer-2}{}\fi
\noindent\textit{Source reference.} Original reference 0B74: sources/stacks-spaces-chow/pinned-original-source.tex, source lines 4017--4017. Prerequisite pointer invoked at proof line(s) 4530; it names a standard fact of the same chapter that the counted argument uses, and every step of the counted proof built on it is retained verbatim in the delivered body.\par\medskip
\ifdefined\hypertarget\hypertarget{sourcepointer-3}{}\fi
\noindent\textit{Source reference.} Original reference 02TT: sources/stacks-spaces-chow/pinned-original-source.tex, source lines 3948--3948. Prerequisite pointer invoked at proof line(s) 4531; it names a standard fact of the same chapter that the counted argument uses, and every step of the counted proof built on it is retained verbatim in the delivered body.\par\medskip
\ifdefined\hypertarget\hypertarget{sourcepointer-4}{}\fi
\noindent\textit{Source reference.} Original reference 02SU: sources/stacks-spaces-chow/pinned-original-source.tex, source lines 2933--2933. Prerequisite pointer invoked at proof line(s) 4565; it names a standard fact of the same chapter that the counted argument uses, and every step of the counted proof built on it is retained verbatim in the delivered body.\par\medskip
\ifdefined\hypertarget\hypertarget{sourcepointer-5}{}\fi
\noindent\textit{Source reference.} Original reference 02TR: sources/stacks-spaces-chow/pinned-original-source.tex, source lines 3889--3889. Prerequisite pointer invoked at proof line(s) 4564; it names a standard fact of the same chapter that the counted argument uses, and every step of the counted proof built on it is retained verbatim in the delivered body.\par\medskip
\ifdefined\hypertarget\hypertarget{sourcepointer-6}{}\fi
\noindent\textit{Source reference.} Original reference 02RX: sources/stacks-spaces-chow/pinned-original-source.tex, source lines 1951--1951. Prerequisite pointer invoked at proof line(s) 4556; it names a standard fact of the same chapter that the counted argument uses, and every step of the counted proof built on it is retained verbatim in the delivered body.\par\medskip
\ifdefined\hypertarget\hypertarget{sourcepointer-7}{}\fi
\noindent\textit{Source reference.} Original reference 02QL: sources/stacks-spaces-chow/pinned-original-source.tex, source lines 264--264. Prerequisite pointer invoked at proof line(s) 4516; it names a standard fact of the same chapter that the counted argument uses, and every step of the counted proof built on it is retained verbatim in the delivered body.\par\medskip
\ifdefined\hypertarget\hypertarget{sourcepointer-8}{}\fi
\noindent\textit{Source reference.} Original lemma 0847: sources/stacks-chow/tags.txt, source lines 10312--10312. Prerequisite pointer invoked at proof line(s) 4544; it names a standard fact of a sibling chapter that the counted argument uses, and every step of the counted proof built on it is retained verbatim in the delivered body.\par\medskip
\begin{lemma}
\label{lemma-vectorbundle}
In Situation \ref{situation-setup} let $X/B$ be good.
Let $\mathcal{E}$ be a finite locally free sheaf of rank $r$ on $X$.
Let
$$
p :
E = \underline{\Spec}(\text{Sym}^*(\mathcal{E}))
\longrightarrow
X
$$
be the associated vector bundle over $X$.
Then $p^* : \CH_k(X) \to \CH_{k + r}(E)$ is an isomorphism for all $k$.
\end{lemma}

\begin{proof}
(For the case of linebundles, see Lemma \ref{lemma-linebundle}.)
For surjectivity see Lemma \ref{lemma-pullback-affine-fibres-surjective}.
Let $(\pi  : P \to X, \mathcal{O}_P(1))$
be the projective space bundle associated
to the finite locally free sheaf $\mathcal{E} \oplus \mathcal{O}_X$.
Let $s \in \Gamma(P, \mathcal{O}_P(1))$ correspond to the global
section $(0, 1) \in \Gamma(X, \mathcal{E} \oplus \mathcal{O}_X)$.
Let $D = Z(s) \subset P$. Note that
$(\pi|_D : D \to X , \mathcal{O}_P(1)|_D)$
is the projective space bundle associated
to $\mathcal{E}$. We denote $\pi_D = \pi|_D$ and
$\mathcal{O}_D(1) = \mathcal{O}_P(1)|_D$.
Moreover, $D$ is an effective
Cartier divisor on $P$. Hence $\mathcal{O}_P(D) = \mathcal{O}_P(1)$
(see Divisors on Spaces, Lemma \ref{spaces-divisors-lemma-characterize-OD}).
Also there is an isomorphism
$E \cong P \setminus D$. Denote $j : E \to P$ the
corresponding open immersion.
For injectivity we use that the kernel of
$$
j^* :
\CH_{k + r}(P)
\longrightarrow
\CH_{k + r}(E)
$$
are the cycles supported in the effective Cartier divisor $D$,
see Lemma \ref{lemma-restrict-to-open}. So if $p^*\alpha = 0$, then
$\pi^*\alpha = i_*\beta$ for some $\beta \in \CH_{k + r}(D)$.
By Lemma \ref{lemma-chow-ring-projective-bundle} we may write
$$
\beta = \pi_D^*\beta_0 +
\ldots + c_1(\mathcal{O}_D(1))^{r - 1} \cap \pi_D^* \beta_{r - 1}.
$$
for some $\beta_i \in \CH_{k + i}(X)$.
By Lemmas \ref{lemma-relative-effective-cartier}
and \ref{lemma-pushforward-cap-c1}
this implies
$$
\pi^*\alpha = i_*\beta =
c_1(\mathcal{O}_P(1)) \cap \pi^*\beta_0 +
\ldots +
c_1(\mathcal{O}_D(1))^r \cap \pi^*\beta_{r - 1}.
$$
Since the rank of $\mathcal{E} \oplus \mathcal{O}_X$ is $r + 1$
this contradicts Lemma \ref{lemma-pushforward-cap-c1} unless all
$\alpha$ and all $\beta_i$ are zero.
\end{proof}

\end{document}
